% SGA 3, Expose II, Loop 1 English reconstruction.
% Scope: Section 4.1 through the end of Example 4.4.2, stopping immediately
% before Proposition 4.5.
% Authority: Polo--Gille Exp2-14oct24.pdf, component-local pages 22--34
% (combined-reader pages 80--92; printed pages 70--82). Force-OCR is a
% drafting locator; English wording was compared with Jacob Reinhold's
% CC BY 4.0 Markdown lineage.

\setcounter{footnote}{48}

% Source page: Exp2 component-local p. 22; combined-reader p. 80; printed p. 70.
\section*{4. Tangent Space to a Group --- Lie Algebras}
\addcontentsline{toc}{section}{4. Tangent Space to a Group --- Lie Algebras}

\subsection*{4.1.}
Let \(G\) be a group functor over \(S\). By 3.9.0.2,
\(T_{G/S}(\mathcal M)\) and
\(\underline{\operatorname{Lie}}(G/S,\mathcal M)\) are endowed with
group structures over \(S\), and there are group morphisms

\SGAThreeDiagram[0.52\textwidth]
  {figures/exp2/exp2_p022_split_tangent_sequence.png}
  {The split tangent sequence of the group functor \(G\).}
  {fig:exp2-split-tangent-sequence}
  {Polo--Gille Exp2-14oct24.pdf, component-local p. 22; combined-reader p. 80; printed p. 70}

By definition, \(i\) identifies
\(\underline{\operatorname{Lie}}(G/S,\mathcal M)\) with the kernel of
\(p\), and \(s\) is a section of \(p\). It follows from I~2.3.7 that this
sequence identifies \(T_{G/S}(\mathcal M)\) with the semidirect product
of \(G\) by \(\underline{\operatorname{Lie}}(G/S,\mathcal M)\).

\subsection*{Definition 4.1.A.}
The corresponding action of \(G\) on
\(\underline{\operatorname{Lie}}(G/S,\mathcal M)\) is denoted
\[
  \operatorname{Ad}:G\longrightarrow
  \underline{\operatorname{Aut}}_{S\text{-gr.}}
    \bigl(\underline{\operatorname{Lie}}(G/S,\mathcal M)\bigr)
\]
and is called the \emph{adjoint representation} of \(G\), relative to
\(\mathcal M\). Thus, by definition, for \(x\in G(S')\) and
\(X\in\underline{\operatorname{Lie}}(G/S,\mathcal M)(S')\),
\[
  \operatorname{Ad}(x)X
  =i^{-1}\bigl(s(x)i(X)s(x)^{-1}\bigr).
\]
If \(G\) is commutative, then \(T_{G/S}(\mathcal M)\) is commutative as
well and \(\operatorname{Ad}(x)X=X\).
\footnote{Editor's note: the numbering 4.1.A and 4.1.B has been added in
order to bring these definitions into relief.}

\subsection*{Definition 4.1.B.}
Let \(G,H\) be group functors over \(S\), and let \(f:G\to H\) be a group
morphism. Functoriality gives a morphism of exact sequences compatible
with the canonical sections.

\SGAThreeDiagram[0.66\textwidth]
  {figures/exp2/exp2_p022_derived_exact_sequences.png}
  {Functoriality of the split tangent sequence.}
  {fig:exp2-derived-exact-sequences}
  {Polo--Gille Exp2-14oct24.pdf, component-local p. 22; combined-reader p. 80; printed p. 70}

The morphism \(L(f)\), also denoted
\(\underline{\operatorname{Lie}}(f)\), is called the \emph{derived
morphism} of \(f\).
\footnote{Editor's note: if \(f\) is a monomorphism, then so are \(L(f)\)
and \(T(f)\); cf. 3.7.}

\subsection*{Remark 4.1.C.}
If \(G/S\) and \(H/S\) satisfy condition \(E\), then \(L(f)\) respects the
\(\mathcal O_S\)-module structures obtained from ``functoriality in
\(\mathcal M\)'' (cf. 3.6).
\footnote{Editor's note: this remark has been added.}

\subsection*{Proposition 4.1.1.}
Let \(g\in G(S)\). Then
\[
  \operatorname{Ad}(g):
  \underline{\operatorname{Lie}}(G/S,\mathcal M)
  \longrightarrow
  \underline{\operatorname{Lie}}(G/S,\mathcal M)
\]
is the derived morphism of
\(\operatorname{Int}(g):G\to G\).

Indeed,
\(\operatorname{Ad}(g)X=i^{-1}(\operatorname{Int}(g)i(X))\), which is
precisely \(L(\operatorname{Int}(g))(X)\) by the definition of the
derived morphism.

Suppose that \(G/S\) satisfies \(E\). By Proposition~3.9, the group
structure just defined on
\(\underline{\operatorname{Lie}}(G/S,\mathcal M)\) is exactly the one
induced by its \(\mathcal O_S\)-module structure (which is defined by
(E)). Proposition~4.1.1 and the functoriality of the
\(\mathcal O_S\)-action (cf. 3.6) therefore yield the following
corollary.

% Source page: Exp2 component-local p. 23; combined-reader p. 81; printed p. 71.
\subsection*{Corollary 4.1.1.1.}
Suppose that \(G/S\) satisfies \(E\). Then \(\operatorname{Ad}\) maps \(G\)
into the subgroup
\[
  \underline{\operatorname{Aut}}_{\mathcal O_S\text{-mod.}}
    \bigl(\underline{\operatorname{Lie}}(G/S,\mathcal M)\bigr)
\]
of
\(\underline{\operatorname{Aut}}_{S\text{-gr.}}
(\underline{\operatorname{Lie}}(G/S,\mathcal M))\). In other words, for
every \(g\in G(S')\), the automorphism \(\operatorname{Ad}(g)\) respects
the \(\mathcal O_{S'}\)-module structure on
\(\underline{\operatorname{Lie}}(G'/S',\mathcal M)\), where
\(G'=G\times_SS'\). Thus \(\operatorname{Ad}\) is a \emph{linear
representation} (cf. I~3.2) of \(G\) on the \(\mathcal O_S\)-module
\(\underline{\operatorname{Lie}}(G/S,\mathcal M)\).

\subsection*{Remarks 4.1.1.2.}
\textit{a)} The functor \(G/S\) satisfies \(E\) if and only if, for every
pair \((\mathcal M,\mathcal N)\) of finite free
\(\mathcal O_S\)-modules, the following diagram, obtained by applying
\(\underline{\operatorname{Lie}}(G/S,-)\) to diagram \((*)\) of 2.1, is
cartesian.

\SGAThreeDiagram[0.45\textwidth]
  {figures/exp2/exp2_p023_condition_e_cartesian.png}
  {The cartesian criterion for condition \(E\).}
  {fig:exp2-condition-e-cartesian}
  {Polo--Gille Exp2-14oct24.pdf, component-local p. 23; combined-reader p. 81; printed p. 71}

\textit{b)} Suppose that \(G/S\) satisfies \(E\). The derived morphism of
the group law
\(\pi:G\times_SG\to G\) is the addition law on
\(\underline{\operatorname{Lie}}(G/S,\mathcal M)\). Notice that
\(\pi\) is not a group morphism, but \(\pi(e,e)=e\); hence the derived
morphism \(L(\pi)\) maps
\[
  T^{(e,e)}_{(G\times_SG)/S}(\mathcal M)
  =\underline{\operatorname{Lie}}(G/S,\mathcal M)
   \times_S\underline{\operatorname{Lie}}(G/S,\mathcal M)
\]
to \(\underline{\operatorname{Lie}}(G/S,\mathcal M)\); cf. 3.7 and
3.8. Likewise, for \(n\in\mathbb Z\), let \(n_G:G\to G\) denote the
\(S\)-functor morphism \(g\mapsto g^n\). Then \(L(n_G)\) is
multiplication by \(n\) on \(\underline{\operatorname{Lie}}(G/S)\); cf.
3.9.4.

\subsection*{4.1.2.0.}
Consider the \(S\)-functor
\(
  \underline{\operatorname{Hom}}_{G/S}
    (G,T_{G/S}(\mathcal M)).
\)\footnote{Editor's note: the original has been expanded in what follows
to display the action of \(G\) on the functors
\(\underline{\operatorname{Hom}}_{G/S}(G,T_{G/S}(\mathcal M))\) and
\(\underline{\operatorname{Hom}}_S
(G,\underline{\operatorname{Lie}}(G/S,\mathcal M))\).}
For every \(S'\to S\), one has
\(T_{G/S}(\mathcal M)_{S'}\simeq T_{G_{S'}/S'}(\mathcal M)\)
(cf. 3.4), and therefore
\[
\begin{aligned}
 &\underline{\operatorname{Hom}}_{G/S}
   (G,T_{G/S}(\mathcal M))(S')\\
 &\qquad\simeq
 \operatorname{Hom}_{G_{S'}}
   \bigl(G_{S'},T_{G_{S'}/S'}(\mathcal M)\bigr)
 =\Gamma\bigl(T_{G_{S'}/S'}(\mathcal M)/G_{S'}\bigr).
\end{aligned}
\]
First note the isomorphism, functorial in \(S'\),
\begin{equation}
 \operatorname{Hom}_{S'}
  \bigl(G_{S'},\underline{\operatorname{Lie}}(G_{S'}/S',\mathcal M)\bigr)
 \xrightarrow{\sim}
 \Gamma\bigl(T_{G_{S'}/S'}(\mathcal M)/G_{S'}\bigr),
 \tag{*}
\end{equation}
which sends \(f:G_{S'}\to
\underline{\operatorname{Lie}}(G_{S'}/S',\mathcal M)\) to the section
\(s_f:G_{S'}\to T_{G_{S'}/S'}(\mathcal M)\) defined, for every
\(S''\to S'\) and \(g\in G(S'')\), by
\[
  s_f(g)=i(f(g))s(g).
\]

% Source page: Exp2 component-local p. 24; combined-reader p. 82; printed p. 72.
Let \(h\) be an automorphism of the functor \(G_{S'}\) over \(S'\), not
necessarily preserving the group structure. To every section \(\tau\) of
\(T_{G_{S'}/S'}(\mathcal M)\) one associates, by transport of structure,
a section \(h(\tau)\): it is the unique section that makes the following
diagram commute.

\SGAThreeDiagram[0.43\textwidth]
  {figures/exp2/exp2_p024_transport_square.png}
  {Transport of a tangent section by an automorphism \(h\).}
  {fig:exp2-transport-square}
  {Polo--Gille Exp2-14oct24.pdf, component-local p. 24; combined-reader p. 82; printed p. 72}

In particular, take \(h=t_x\), right translation by an element
\(x\in G(S')\); thus \(h(g)=t_x(g)=g\cdot x\) for all
\(g\in G(S'')\), \(S''\to S'\). One then has
\[
  t_x(s_f)=s_{t_x(f)},
\]
where \(t_x(f):G_{S'}\to
\underline{\operatorname{Lie}}(G_{S'}/S',\mathcal M)\) is given by
\[
  t_x(f)(g)=f(g\cdot x^{-1}).
\]

Let \(G\) act by right translations on
\[
 \underline{\operatorname{Hom}}_{G/S}(G,T_{G/S}(\mathcal M))
 \quad\text{and}\quad
 \underline{\operatorname{Hom}}_S
   (G,\underline{\operatorname{Lie}}(G/S,\mathcal M))
\]
as follows. For \(S'\to S\) and \(x\in G(S')\), let
\[
\begin{aligned}
 \sigma&\in\Gamma(T_{G_{S'}/S'}(\mathcal M)/G_{S'}),\\
 f&\in\operatorname{Hom}_{S'}
 (G_{S'},\underline{\operatorname{Lie}}(G_{S'}/S',\mathcal M)).
\end{aligned}
\]
Set
\[
 (\sigma\cdot x)(g)=\sigma(g\cdot x^{-1})\cdot s(x),
 \qquad
 (f\cdot x)(g)=f(g\cdot x^{-1})
\]
for every \(g\in G(S'')\), \(S''\to S'\). The isomorphism \((*)\)
respects these \(G\)-actions.

Consequently, under this isomorphism, the elements of
\[
 \underline{\operatorname{Hom}}_{G/S}
   (G,T_{G/S}(\mathcal M))^G(S')
\]
correspond to the constant morphisms from \(G_{S'}\) to
\(\underline{\operatorname{Lie}}(G_{S'}/S',\mathcal M)\), that is, to
the morphisms factoring through \(G_{S'}\to S'\); equivalently, they
correspond to the elements of
\(\underline{\operatorname{Lie}}(G/S,\mathcal M)(S')\).

\paragraph{Terminology.}
The elements of
\(\underline{\operatorname{Hom}}_{G/S}
(G,T_{G/S}(\mathcal M))^G(S')\) are called the sections of
\(T_{G_{S'}/S'}(\mathcal M)\) invariant under right translation.

\subsection*{Proposition 4.1.2.}
The map
\[
 \underline{\operatorname{Lie}}(G/S,\mathcal M)(S)
 \longrightarrow \Gamma(T_{G/S}(\mathcal M)/G),
 \qquad X\longmapsto (x\mapsto X\cdot x),
\]
is a bijection onto the set of sections invariant under right
translation.
\footnote{Statements 4.1.2, 4.1.3, and 4.1.4 can be obtained more simply
by observing that the automorphisms of \(G\) invariant under right
translations are the left translations. Editor's note: see, for example,
the proofs of Propositions 4.6 and 2.5 of \textup{[DG70]}, \S II.4.}

% Source page: Exp2 component-local p. 25; combined-reader p. 83; printed p. 73.
Likewise, let \(G\) act on the right on
\(\underline{\operatorname{End}}_{I_S(\mathcal M)/S}
(G_{I_S(\mathcal M)})\) by
\[
 (u\cdot x)(g)=u(g\cdot x^{-1})\cdot x,
\]
for \(S'\to S\), \(x\in G(S')\),
\(u\in\operatorname{End}_{I_{S'}(\mathcal M)}
(G_{I_{S'}(\mathcal M)})\), and \(g\in G(S'')\),
\(S''\to I_{S'}(\mathcal M)\). The morphism of 3.12.1
\[
 \underline{\operatorname{Hom}}_{G/S}(G,T_{G/S}(\mathcal M))
 \longrightarrow
 \underline{\operatorname{End}}_{I_S(\mathcal M)/S}
   (G_{I_S(\mathcal M)})
\]
respects the \(G\)-actions. For each \(S'\to S\), it therefore induces a
bijection between
\(\Gamma(T_{G_{S'}/S'}(\mathcal M)/G_{S'})\) and the set of
\(I_{S'}(\mathcal M)\)-endomorphisms \(u\) of
\(G_{I_{S'}(\mathcal M)}\) that induce the identity on \(G\) and
``commute with right translations,'' meaning that
\(u_{S''}\cdot x=u_{S''}\) for every \(S''\to S'\) and
\(x\in G(S'')\).

\subsection*{Proposition 4.1.3.}
There is a bijection, functorial in \(G\), between
\(\underline{\operatorname{Lie}}(G/S,\mathcal M)(S)\) and the set of
\(I_S(\mathcal M)\)-endomorphisms of \(G_{I_S(\mathcal M)}\) that
induce the identity on \(G\) and commute with the right translations of
\(G\), in the sense just specified.

Taking 3.13 into account gives the following theorem.

\subsection*{Theorem 4.1.4.}
Let \(G\) be an \(S\)-functor in groups, and suppose that \(G/S\)
satisfies \(E\). Then
\(\underline{\operatorname{Lie}}(G/S,\mathcal M)(S)\) is identified,
functorially in \(G\), with the group of
\(I_S(\mathcal M)\)-automorphisms of \(G_{I_S(\mathcal M)}\) that
induce the identity on \(G\) and commute with the right translations of
\(G\).

For \(\mathcal M=\mathcal O_S\), this recovers one of the classical
definitions of the Lie algebra of a group.

\subsection*{4.2.0.}
Before proceeding, let us establish further consequences of
3.11.\footnote{Editor's note: paragraph 4.2.0 has been added; its results
are used implicitly in 4.2 and 4.7 of the original.} Let
\(X,Y\) be over \(Z\), with \(Z\) over \(S\), as in Section~1. As was
seen in 3.11, the isomorphisms 1.1(4)

\SGAThreeDiagram[0.78\textwidth]
  {figures/exp2/exp2_p025_hom_adjunction_triangle.png}
  {The adjunction isomorphisms of 1.1(4).}
  {fig:exp2-hom-adjunction-p25}
  {Polo--Gille Exp2-14oct24.pdf, component-local p. 25; combined-reader p. 83; printed p. 73}

induce the following isomorphism \(\theta\).

\SGAThreeDiagram[0.78\textwidth]
  {figures/exp2/exp2_p025_tangent_adjunction_triangle.png}
  {The induced tangent-space isomorphism \(\theta\).}
  {fig:exp2-tangent-adjunction-p25}
  {Polo--Gille Exp2-14oct24.pdf, component-local p. 25; combined-reader p. 83; printed p. 73}

% Source page: Exp2 component-local p. 26; combined-reader p. 84; printed p. 74.
By 1.3, if \(Y\) is a \(Z\)-group, then
\(\underline{\operatorname{Hom}}_Z(V,Y)\) is a \(Z\)-group for every
\(V\to Z\), in particular for \(V=I_Z(\mathcal M)\). Explicitly, if
\(Z''\to Z'\to Z\) and
\(\varphi,\psi\in\operatorname{Hom}_Z(V_{Z'},Y)(Z')\), their product is
defined by
\[
  (\varphi\cdot\psi)(v)=\varphi(v)\psi(v)
  \qquad (v\in V_{Z'}(Z'')).
\]

Suppose now that \(X\) and \(Y\) are \(Z\)-groups.

\subsection*{Definition 4.2.0.1.}
We denote by
\(\underline{\operatorname{Hom}}_{(Z/S)\text{-gr.}}(X,Y)\) the
subfunctor of \(\underline{\operatorname{Hom}}_{Z/S}(X,Y)\) defined, for
every \(S'\to S\), by
\begin{equation}
 \underline{\operatorname{Hom}}_{(Z/S)\text{-gr.}}(X,Y)(S')
 =\operatorname{Hom}_{Z_{S'}\text{-gr.}}(X_{S'},Y_{S'}).
 \tag{3}
\end{equation}
The same definition applies with \(Y\) replaced by the \(Z\)-group
\(T_{Y/Z}(\mathcal M)\).

Under the isomorphisms (2) above,
\(T_{\underline{\operatorname{Hom}}_{(Z/S)\text{-gr.}}(X,Y)/S}
(\mathcal M)(S')\) corresponds to the \(Z_{S'}\)-morphisms
\[
 \varphi:X_{S'}\times_{S'}I_{S'}(\mathcal M)\longrightarrow Y_{S'}
\]
that are multiplicative in \(X\), that is,
\[
 \varphi(x_1x_2,m)=\varphi(x_1,m)\varphi(x_2,m).
\]
These in turn correspond to the morphisms of \(Z_{S'}\)-groups
\(X_{S'}\to T_{Y/Z}(\mathcal M)_{S'}\). Hence:

\subsection*{Proposition 4.2.0.2.}
Let \(X,Y\) be \(Z\)-groups, where \(Z\) is over \(S\). There are
isomorphisms of \(S\)-functors, functorial in \(\mathcal M\),
\[
 T_{\underline{\operatorname{Hom}}_{(Z/S)\text{-gr.}}(X,Y)/S}
   (\mathcal M)
 \xrightarrow{\sim}
 \underline{\operatorname{Hom}}_{(Z/S)\text{-gr.}}
   (X,T_{Y/Z}(\mathcal M)).
\]

Taking \(Z=S\) gives the following consequence. If \(Y\) is a
commutative \(S\)-group, then so is \(T_{Y/S}(\mathcal M)\), and hence
so are
\(H=\underline{\operatorname{Hom}}_{S\text{-gr.}}(X,Y)\),
\(\underline{\operatorname{Hom}}_{S\text{-gr.}}
(X,T_{Y/S}(\mathcal M))\), and \(T_{H/S}(\mathcal M)\).

\subsection*{Corollary 4.2.0.3.}
Let \(X,Y\) be \(S\)-groups. There are isomorphisms of \(S\)-functors,
functorial in \(\mathcal M\),
\[
 T_{\underline{\operatorname{Hom}}_{S\text{-gr.}}(X,Y)/S}
   (\mathcal M)
 \xrightarrow{\sim}
 \underline{\operatorname{Hom}}_{S\text{-gr.}}
   (X,T_{Y/S}(\mathcal M)).
\]
If \(Y\) is commutative, these are moreover isomorphisms of abelian
\(S\)-groups.

\subsection*{Definition 4.2.0.4.}
If \(Y\) is an \(\mathcal O_S\)-module, the functor
\(T_{Y/S}(\mathcal M)\), respectively
\(\underline{\operatorname{Lie}}(Y/S,\mathcal M)\), carries an
\(\mathcal O_S\)-module structure induced by that of \(Y\). Equipped
with this structure, it will be denoted
\(T'_{Y/S}(\mathcal M)\), respectively
\(\underline{\operatorname{Lie}}'(Y/S,\mathcal M)\).
\footnote{Editor's note: the prime is used to distinguish the
\(\mathcal O_S\)-module structure induced by that of \(Y\) from the one
that will be obtained from condition \(E\).}

Consequently, if \(X,Y\) are \(\mathcal O_S\)-modules, then
\(T'_{Y/S}(\mathcal M)=
\underline{\operatorname{Hom}}_S(I_S(\mathcal M),Y)\). If
\(H=\underline{\operatorname{Hom}}_{\mathcal O_S\text{-mod.}}(X,Y)\),
then
\(\underline{\operatorname{Hom}}_{\mathcal O_S\text{-mod.}}
(X,T'_{Y/S}(\mathcal M))\) and \(T'_{H/S}(\mathcal M)\) also carry
\(\mathcal O_S\)-module structures.

% Source page: Exp2 component-local p. 27; combined-reader p. 85; printed p. 75.
\subsection*{Corollary 4.2.0.5.}
If \(X,Y\) are \(\mathcal O_S\)-modules, there are isomorphisms of
\(\mathcal O_S\)-modules, functorial in \(\mathcal M\),
\[
 T'_{\underline{\operatorname{Hom}}_{\mathcal O_S\text{-mod.}}(X,Y)/S}
   (\mathcal M)
 \xrightarrow{\sim}
 \underline{\operatorname{Hom}}_{\mathcal O_S\text{-mod.}}
   (X,T'_{Y/S}(\mathcal M)).
\]

\subsection*{Definition 4.2.A.}
Let \(X,L\) be \(S\)-groups, with \(X\) acting on \(L\) by group
automorphisms (cf. I~2.3.5). We define the subfunctor
\(\underline Z^1_S(X,L)\) of
\(\underline{\operatorname{Hom}}_S(X,L)\) by
\[
\begin{aligned}
 \underline Z^1_S(X,L)(S')
 =\{\,\varphi\in\operatorname{Hom}_{S'}(X_{S'},L_{S'})\mid{}
 &\varphi(x_1x_2)
   =\varphi(x_1)\bigl(x_1\cdot\varphi(x_2)\bigr),\\[-2pt]
 &x_1,x_2\in X(S''),\ S''\to S'\,\}.
\end{aligned}
\]
It is called the \emph{functor of crossed homomorphisms} from \(X\) to
\(L\).
\footnote{Editor's note: the original has been expanded below; in
particular, Definition 4.2.A and Remark 4.2.B have been added.}

\subsection*{Remark 4.2.B.}
\textit{a)} If \(L'\) is another \(S\)-group on which \(X\) acts by group
automorphisms, then
\[
 \underline Z^1_S(X,L\times_SL')
 \simeq \underline Z^1_S(X,L)\times_S\underline Z^1_S(X,L').
\]

\textit{b)} If \(L\) is an \(X\)-\(\mathcal O_S\)-module, then
\(\underline Z^1_S(X,L)\) is the kernel of the differential
\[
 \partial:\underline{\operatorname{Hom}}_S(X,L)
 \longrightarrow\underline{\operatorname{Hom}}_S(X^2,L)
\]
defined in I~5.1. In particular, \(\underline Z^1_S(X,L)\) is then an
\(\mathcal O_S\)-module.

\subsection*{4.2.}
Let \(u:X\to Y\) be a morphism of \(S\)-groups. We describe
\(L^u_{\underline{\operatorname{Hom}}_{S\text{-gr.}}(X,Y)/S}
(\mathcal M)\). First, 3.11.3 gives an isomorphism of \(S\)-functors,
functorial in \(\mathcal M\),
\begin{equation}
 L^u_{\underline{\operatorname{Hom}}_S(X,Y)/S}(\mathcal M)
 \xrightarrow{\sim}
 \underline{\operatorname{Hom}}_{Y/S}(X,T_{Y/S}(\mathcal M)).
 \tag{\(\dagger\)}
\end{equation}
Since \(Y\) is an \(S\)-group,
\[
 T_{Y/S}(\mathcal M)
 =\underline{\operatorname{Lie}}(Y/S,\mathcal M)\cdot Y
 =\underline{\operatorname{Lie}}(Y/S,\mathcal M)_Y,
\]
and one obtains an isomorphism, functorial in \(\mathcal M\),
\[
 L^u_{\underline{\operatorname{Hom}}_S(X,Y)/S}(\mathcal M)
 \xrightarrow{\sim}
 \underline{\operatorname{Hom}}_S
   (X,\underline{\operatorname{Lie}}(Y/S,\mathcal M)).
\]

For \(S'\to S\), let \(u':X'\to Y'\) be the base change of \(u\). The
\(S\)-functor\footnote{Editor's note: this is Definition 4.2.0.1 applied
with \(Z=Y\), to the \(Y\)-groups \(u:X\to Y\) and
\(T_{Y/S}(\mathcal M)\to Y\).}
\[
 \underline{\operatorname{Hom}}_{(Y/S)\text{-gr.}}
   \bigl(X,\underline{\operatorname{Lie}}(Y/S,\mathcal M)\cdot Y\bigr)
\]
is defined by
\[
\begin{aligned}
 &\underline{\operatorname{Hom}}_{(Y/S)\text{-gr.}}
 \bigl(X,\underline{\operatorname{Lie}}(Y/S,\mathcal M)\cdot Y\bigr)(S')\\
 &\quad=
 \operatorname{Hom}_{Y'\text{-gr.}}
 \bigl(X',(\underline{\operatorname{Lie}}(Y/S,\mathcal M)\cdot Y)_{S'}\bigr)\\
 &\quad=
 \operatorname{Hom}_{Y'\text{-gr.}}
 \bigl(X',\underline{\operatorname{Lie}}(Y'/S',\mathcal M)\cdot Y'\bigr).
\end{aligned}
\]
Thus \((\dagger)\) induces an isomorphism
\begin{equation}
 L^u_{\underline{\operatorname{Hom}}_{S\text{-gr.}}(X,Y)/S}
   (\mathcal M)
 \xrightarrow{\sim}
 \underline{\operatorname{Hom}}_{(Y/S)\text{-gr.}}
 \bigl(X,\underline{\operatorname{Lie}}(Y/S,\mathcal M)\cdot Y\bigr).
 \tag{\(\dagger'\)}
\end{equation}

On the other hand, the composite

\SGAThreeDiagram[0.50\textwidth]
  {figures/exp2/exp2_p027_adjoint_action_chain.png}
  {The action of \(X\) on
  \(\underline{\operatorname{Lie}}(Y/S,\mathcal M)\).}
  {fig:exp2-adjoint-action-chain}
  {Polo--Gille Exp2-14oct24.pdf, component-local p. 27; combined-reader p. 85; printed p. 75}

defines an action of \(X\) on
\(L=\underline{\operatorname{Lie}}(Y/S,\mathcal M)\) by group
automorphisms.

% Source page: Exp2 component-local p. 28; combined-reader p. 86; printed p. 76.
Let
\(\Phi\in\underline{\operatorname{Hom}}_{Y/S}
(X,\underline{\operatorname{Lie}}(Y/S,\mathcal M)\cdot Y)\). For
\(S''\to S'\to S\) and \(x\in X(S'')\), there is a unique expression
\[
 \Phi(S')(x)=\varphi(S')(x)\cdot u'(x),
 \qquad
 \varphi(S')(x)\in
 \underline{\operatorname{Lie}}(Y'/S',\mathcal M)(S'').
\]
This determines an element
\(\varphi\in\underline{\operatorname{Hom}}_S
(X,\underline{\operatorname{Lie}}(Y/S,\mathcal M))\). The map
\(\Phi(S')\) is a group morphism if and only if, for all
\(x_1,x_2\in X(S'')\),
\[
\begin{aligned}
 \varphi(S')(x_1x_2)
 &=\varphi(S')(x_1)
   \bigl(u(x_1)\varphi(S')(x_2)u(x_1)^{-1}\bigr)\\
 &=\varphi(S')(x_1)
   \bigl(x_1\cdot\varphi(S')(x_2)\bigr),
\end{aligned}
\]
that is, precisely when
\(\varphi\in\underline Z^1_S
(X,\underline{\operatorname{Lie}}(Y/S,\mathcal M))\).

\subsection*{Proposition 4.2.}
Let \(u:X\to Y\) be a morphism of \(S\)-groups. There is an isomorphism
of \(S\)-functors, functorial in \(\mathcal M\),
\[
 L^u_{\underline{\operatorname{Hom}}_{S\text{-gr.}}(X,Y)/S}
   (\mathcal M)
 \xrightarrow{\sim}
 \underline Z^1_S
   \bigl(X,\underline{\operatorname{Lie}}(Y/S,\mathcal M)\bigr).
\]

Suppose in addition that \(Y/S\) satisfies \(E\). It follows from
4.2.0.3, just as in the proof of 3.11.1, that
\(\underline{\operatorname{Hom}}_{S\text{-gr.}}(X,Y)/S\) satisfies
\(E\). Thus (cf. 3.5.1)
\[
\begin{aligned}
 L^u_{\underline{\operatorname{Hom}}_{S\text{-gr.}}(X,Y)/S}
 (\mathcal M\oplus\mathcal N)
 \simeq{}&
 L^u_{\underline{\operatorname{Hom}}_{S\text{-gr.}}(X,Y)/S}
 (\mathcal M)\\
 &\times_S
 L^u_{\underline{\operatorname{Hom}}_{S\text{-gr.}}(X,Y)/S}
 (\mathcal N).
\end{aligned}
\]
This also follows directly from Proposition~4.2 and Remark~4.2.B\(a\).
Consequently, both sides of the isomorphism in Proposition~4.2 carry
\(\mathcal O_S\)-module structures obtained from functoriality in
\(\mathcal M\), and that isomorphism respects them.

\subsection*{Proposition 4.2 bis.}
Let \(u:X\to Y\) be a morphism of \(S\)-groups, and suppose that \(Y/S\)
satisfies \(E\). There is an isomorphism of \(\mathcal O_S\)-modules,
functorial in \(\mathcal M\),
\[
 L^u_{\underline{\operatorname{Hom}}_{S\text{-gr.}}(X,Y)/S}
   (\mathcal M)
 \xrightarrow{\sim}
 \underline Z^1_S
   \bigl(X,\underline{\operatorname{Lie}}(Y/S,\mathcal M)\bigr).
\]
\footnote{Editor's note: the original has been expanded from this point;
in particular, Proposition 4.2 bis has been added.}

When \(Y/S\) satisfies \(E\), 3.13.1 also gives, just as in the proof of
3.13.2, an isomorphism functorial in \(\mathcal M\), for every
\(u\in\underline{\operatorname{Isom}}_{S\text{-gr.}}(X,Y)\),
\begin{equation}
 L^u_{\underline{\operatorname{Isom}}_{S\text{-gr.}}(X,Y)/S}
   (\mathcal M)
 \xrightarrow{\sim}
 L^u_{\underline{\operatorname{Hom}}_{S\text{-gr.}}(X,Y)/S}
   (\mathcal M).
 \tag{*}
\end{equation}
This yields the following two corollaries.

\subsection*{Corollary 4.2.1.}
Let \(u:X\to Y\) be a morphism of \(S\)-groups. If \(Y/S\) satisfies
\(E\), there is an isomorphism of \(\mathcal O_S\)-modules, functorial in
\(\mathcal M\),
\[
 L^u_{\underline{\operatorname{Isom}}_{S\text{-gr.}}(X,Y)/S}
   (\mathcal M)
 \xrightarrow{\sim}
 \underline Z^1_S
   \bigl(X,\underline{\operatorname{Lie}}(Y/S,\mathcal M)\bigr).
\]

% Source page: Exp2 component-local p. 29; combined-reader p. 87; printed p. 77.
\subsection*{Corollary 4.2.2.}
Let \(X\) be an \(S\)-group. If \(X/S\) satisfies \(E\), there is an
isomorphism of \(\mathcal O_S\)-modules, functorial in \(\mathcal M\),
\[
 \underline{\operatorname{Lie}}
   \bigl(\underline{\operatorname{Aut}}_{S\text{-gr.}}(X)/S,
         \mathcal M\bigr)
 \xrightarrow{\sim}
 \underline Z^1_S
   \bigl(X,\underline{\operatorname{Lie}}(X/S,\mathcal M)\bigr).
\]

If \(Y\) is commutative, its adjoint action on
\(L=\underline{\operatorname{Lie}}(Y/S,\mathcal M)\) is trivial, so
\(\underline Z^1_S(X,L)=
\underline{\operatorname{Hom}}_{S\text{-gr.}}(X,L)\).
\footnote{Editor's note: the following sentence has been added.}

\subsection*{Corollary 4.2.3.}
Let \(Y\) be a commutative \(S\)-group. There is an isomorphism of
\(S\)-functors, functorial in \(\mathcal M\),
\[
 L^u_{\underline{\operatorname{Hom}}_{S\text{-gr.}}(X,Y)/S}
   (\mathcal M)
 \xrightarrow{\sim}
 \underline{\operatorname{Hom}}_{S\text{-gr.}}
   \bigl(X,\underline{\operatorname{Lie}}(Y/S,\mathcal M)\bigr).
\]

\subsection*{4.3.}
We now consider the case in which \(X\) and \(Y\) are
\(\mathcal O_S\)-modules. Recall (cf. 4.2.0.4) that
\(T'_{Y/S}(\mathcal M)\), respectively
\(\underline{\operatorname{Lie}}'(Y/S,\mathcal M)\), denotes the
functor \(T_{Y/S}(\mathcal M)\), respectively
\(\underline{\operatorname{Lie}}(Y/S,\mathcal M)\), equipped with the
\(\mathcal O_S\)-module structure induced by that of \(Y\).

When \(Y/S\) satisfies \(E\), we shall always write
\(\underline{\operatorname{Lie}}(Y/S,\mathcal M)\) for this functor
equipped instead with the \(\mathcal O_S\)-module structure defined for
every functor satisfying \(E\). By 3.9, the underlying abelian-group
structures of
\(\underline{\operatorname{Lie}}(Y/S,\mathcal M)\) and
\(\underline{\operatorname{Lie}}'(Y/S,\mathcal M)\) coincide, but their
module structures need not coincide (see the counterexample in 6.3). For
\(S'\to S\), \(a\in\mathcal O(S')\), and
\(m\in\underline{\operatorname{Lie}}(Y/S,\mathcal M)(S')\), we write
\(a\mathbin{\cdot'}m\), respectively \(a\cdot m\), for the action coming
from \(Y\), respectively the action obtained from condition \(E\). We use
the same convention for the two actions on \(T_{Y/S}(\mathcal M)\).

There is an isomorphism of \(\mathcal O_S\)-modules
\[
 T'_{Y/S}(\mathcal M)
 \simeq \underline{\operatorname{Lie}}'(Y/S,\mathcal M)\oplus Y.
\]
Exactly as for Propositions~4.2 and 4.2 bis, one therefore obtains:

\subsection*{Proposition 4.3.}
Let \(u:X\to Y\) be a morphism of \(\mathcal O_S\)-modules. There is an
isomorphism of \(S\)-functors, functorial in \(\mathcal M\),
\begin{equation}
 L^u_{\underline{\operatorname{Hom}}_{\mathcal O_S\text{-mod.}}(X,Y)/S}
   (\mathcal M)
 \xrightarrow{\sim}
 \underline{\operatorname{Hom}}_{\mathcal O_S\text{-mod.}}
   \bigl(X,\underline{\operatorname{Lie}}'(Y/S,\mathcal M)\bigr).
 \tag{*}
\end{equation}
If \(Y/S\) satisfies \(E\), then
\(\underline{\operatorname{Hom}}_{\mathcal O_S\text{-mod.}}(X,Y)/S\)
satisfies \(E\), and \((*)\) is an isomorphism of
\(\mathcal O_S\)-modules when both sides carry the structures obtained
from functoriality in \(\mathcal M\).
\footnote{Editor's note: the following sentence has been added.}
\footnote{Editor's note: on the right-hand side, the scalar action is
\((a\varphi)(x)=a\varphi(x)\), where \(a\varphi(x)\) denotes the action
of \(a\) on \(\varphi(x)\in
\underline{\operatorname{Lie}}(Y/S,\mathcal M)\). This differs from the
action \((a\mathbin{\cdot'}\varphi)(x)
=a\mathbin{\cdot'}\varphi(x)=\varphi(ax)\), but the two agree when
\(\underline{\operatorname{Lie}}(Y/S,\mathcal M)
=\underline{\operatorname{Lie}}'(Y/S,\mathcal M)\).}

\subsection*{Remark 4.3 bis.}
Let \(u:X\to Y\) be a morphism of
\(\mathcal O_S\)-modules.\footnote{Editor's note: this remark has been
added; it will be useful in 4.5.1.} Denote by
\(\tau_u\) the map that sends each \(\mathcal O_S\)-linear morphism
\(\varphi:X\to\underline{\operatorname{Lie}}'(Y/S,\mathcal M)\) to
\[
 u\oplus\varphi:X\longrightarrow
 T'_{Y/S}(\mathcal M)
 =Y\oplus\underline{\operatorname{Lie}}'(Y/S,\mathcal M).
 \]

% Source page: Exp2 component-local p. 30; combined-reader p. 88; printed p. 78.
The isomorphism of 4.3 belongs to the following commutative diagram,
functorial in \(\mathcal M\).

\SGAThreeDiagram[0.70\textwidth]
  {figures/exp2/exp2_p030_tau_u_square.png}
  {Compatibility of the isomorphism of 4.3 with \(\tau_u\).}
  {fig:exp2-tau-u-square}
  {Polo--Gille Exp2-14oct24.pdf, component-local p. 30; combined-reader p. 88; printed p. 78}

Moreover, when \(Y/S\) satisfies \(E\), 3.13.1 gives, just as in the proof
of 3.13.2, for every
\(u\in\underline{\operatorname{Isom}}_{\mathcal O_S\text{-mod.}}(X,Y)\),
\begin{equation}
 L^u_{\underline{\operatorname{Isom}}_{\mathcal O_S\text{-mod.}}(X,Y)/S}
   (\mathcal M)
 =
 L^u_{\underline{\operatorname{Hom}}_{\mathcal O_S\text{-mod.}}(X,Y)/S}
   (\mathcal M).
 \tag{*}
\end{equation}

\subsection*{Corollary 4.3.1.}
Let \(X\) be an \(\mathcal O_S\)-module satisfying \(E\) over \(S\). There
is an isomorphism, functorial in \(\mathcal M\),
\[
 \underline{\operatorname{Lie}}
  \bigl(\underline{\operatorname{Aut}}_{\mathcal O_S\text{-mod.}}(X)/S,
        \mathcal M\bigr)
 \xrightarrow{\sim}
 \underline{\operatorname{Hom}}_{\mathcal O_S\text{-mod.}}
  \bigl(X,\underline{\operatorname{Lie}}'(X/S,\mathcal M)\bigr),
\]
and this isomorphism respects the \(\mathcal O_S\)-module structures
obtained from functoriality in \(\mathcal M\). In particular,
\(\underline{\operatorname{Aut}}_{\mathcal O_S\text{-mod.}}(X)/S\)
satisfies \(E\).
\footnote{Editor's note: compare editor's note (61). The final sentence
and its proof have also been added.}

\paragraph{Proof.}
The first assertion follows from \((*)\) and 4.3. For the second, since
\(X/S\) satisfies \(E\), there are isomorphisms of
\(\mathcal O_S\)-modules
\[
 \underline{\operatorname{Lie}}'(X/S,\mathcal M\oplus\mathcal N)
 \simeq
 \underline{\operatorname{Lie}}'(X/S,\mathcal M)
 \times_S
 \underline{\operatorname{Lie}}'(X/S,\mathcal N).
\]
Consequently,
\[
\begin{aligned}
 &\underline{\operatorname{Lie}}
  \bigl(\underline{\operatorname{Aut}}_{\mathcal O_S\text{-mod.}}(X)/S,
  \mathcal M\oplus\mathcal N\bigr)\\
 &\quad\simeq
 \underline{\operatorname{Lie}}
  \bigl(\underline{\operatorname{Aut}}_{\mathcal O_S\text{-mod.}}(X)/S,
  \mathcal M\bigr)
 \times_S
 \underline{\operatorname{Lie}}
  \bigl(\underline{\operatorname{Aut}}_{\mathcal O_S\text{-mod.}}(X)/S,
  \mathcal N\bigr).
\end{aligned}
\]
Remark~4.1.1.2\(a\) now proves the assertion.

\subsection*{4.3.2.}
Before continuing in this direction, let us examine more closely the
relations among \(Y\), \(\underline{\operatorname{Lie}}(Y/S)\), and
\(\underline{\operatorname{Lie}}'(Y/S)\). First,
\begin{equation}
 \underline{\operatorname{Lie}}(\mathcal O_S/S,\mathcal M)
 =\underline{\operatorname{Lie}}'(\mathcal O_S/S,\mathcal M)
 =\mathbf W(\mathcal M),
 \tag{1}
\end{equation}
where \(\mathbf W(\mathcal M)\) is defined in I~4.6. There is therefore a
canonical isomorphism
\begin{equation}
 d:\mathcal O_S\xrightarrow{\sim}
 \underline{\operatorname{Lie}}(\mathcal O_S/S).
 \tag{2}
\end{equation}

Let \(F\) be an \(\mathcal O_S\)-module. For every
\(S_2\to S_1\to S\), there is a dihomomorphism
\begin{equation}
 \left\{
 \begin{aligned}
   F(S_1)&\longrightarrow F(S_2),\\
   \mathcal O(S_1)&\longrightarrow\mathcal O(S_2),
 \end{aligned}
 \right.
 \tag{3}
\end{equation}
and hence a morphism of \(\mathcal O(S_2)\)-modules
\[
 F(S_1)\otimes_{\mathcal O(S_1)}\mathcal O(S_2)
 \longrightarrow F(S_2).
\]
\footnote{Editor's note: \(S'\to S\) has been replaced by
\(S_2\to S_1\to S\), since both \(S_2\) and \(S_1\) must vary below.
The original has also been expanded in what follows.}

% Source page: Exp2 component-local p. 31; combined-reader p. 89; printed p. 79.
In particular, take \(S_1=S'\) and \(S_2=I_{S'}(\mathcal M)\). One
obtains morphisms of \(\mathcal O(S')\)-modules, functorial in
\(\mathcal M\),
\[
 F(S')\otimes_{\mathcal O(S')}
 T_{\mathcal O_S/S}(\mathcal M)(S')
 \longrightarrow T'_{F/S}(\mathcal M)(S').
\]
As \(S'\) varies, these give morphisms of \(\mathcal O_S\)-modules
\begin{equation}
 F\otimes_{\mathcal O_S}T_{\mathcal O_S/S}(\mathcal M)
 \longrightarrow T'_{F/S}(\mathcal M),
 \tag{4}
\end{equation}
functorial in \(\mathcal M\). Compatibility with the tangent-bundle
projections onto their bases then yields morphisms
\begin{equation}
 F\otimes_{\mathcal O_S}
 \underline{\operatorname{Lie}}(\mathcal O_S/S,\mathcal M)
 \longrightarrow
 \underline{\operatorname{Lie}}'(F/S,\mathcal M),
 \tag{5}
\end{equation}
and the following diagram commutes.

\SGAThreeDiagram[0.78\textwidth]
  {figures/exp2/exp2_p031_tangent_exact_rows.png}
  {The morphisms (4) and (5) in the split tangent sequences.}
  {fig:exp2-tangent-exact-rows}
  {Polo--Gille Exp2-14oct24.pdf, component-local p. 31; combined-reader p. 89; printed p. 79}

The morphisms (5) may be regarded as morphisms of abelian \(S\)-groups
\begin{equation}
 F\otimes_{\mathcal O_S}
 \underline{\operatorname{Lie}}(\mathcal O_S/S,\mathcal M)
 \longrightarrow
 \underline{\operatorname{Lie}}(F/S,\mathcal M).
 \tag{6}
\end{equation}
Tensoring \(F\) with the isomorphism
\(d:\mathcal O_S\xrightarrow{\sim}
\underline{\operatorname{Lie}}(\mathcal O_S/S)\) gives, when
\(\mathcal M=\mathcal O_S\), a morphism of abelian \(S\)-groups
\begin{equation}
 F\xrightarrow{\sim}
 F\otimes_{\mathcal O_S}
 \underline{\operatorname{Lie}}(\mathcal O_S/S)
 \longrightarrow \underline{\operatorname{Lie}}(F/S),
 \tag{7}
\end{equation}
also denoted \(d:F\to\underline{\operatorname{Lie}}(F/S)\).

\subsection*{Remark 4.3.3.}
When \(F/S\) satisfies \(E\), the morphisms (6) and (7) need not be
\(\mathcal O_S\)-linear if both sides are equipped with the module
structures obtained from \(\mathcal M\) by condition \(E\).

For example, let \(k\) be a field of characteristic \(p>0\), let
\(S=\operatorname{Spec}(k)\), and let \(F\) be the
\(\mathcal O_S\)-module that assigns to every \(S\)-scheme \(T\) the
set \(F(T)=\Gamma(T,\mathcal O_T)\), with scalars acting through the
\(p\)-th power:
\[
 r\cdot f=r^pf
 \qquad(r\in\mathcal O(T),\ f\in F(T)).
\]
As an \(S\)-functor in groups, \(F\simeq\mathbf G_{a,S}\). Hence \(F\)
satisfies \(E\) and
\(\underline{\operatorname{Lie}}(F/S)\simeq
\underline{\operatorname{Lie}}(\mathbf G_{a,S}/S)\simeq\mathcal O_S\).
For every \(T\to S\), the canonical morphism
\(d:F(T)\to\mathcal O(T)\) is the identity: it respects the abelian-group
structure but not the \(\mathcal O_S\)-module structure.
\footnote{Editor's note: the original asserted that, when \(F/S\)
satisfies \(E\), (6), and hence (7), are morphisms of
\(\mathcal O_S\)-modules. The counterexample in 6.3 contradicts this.
That assertion has been removed, the preceding example inserted, and
Definition 4.4 modified accordingly. This does not affect what follows.}

% Source page: Exp2 component-local p. 32; combined-reader p. 90; printed p. 80.
\subsection*{Remark 4.3.4.}
The morphisms (4) and (5) can be made explicit as follows. The morphism
\[
 \Theta:F\otimes_{\mathcal O_S}T_{\mathcal O_S/S}(\mathcal M)
 \longrightarrow T'_{F/S}(\mathcal M)
 =\underline{\operatorname{Hom}}_S(I_S(\mathcal M),F)
\]
is defined, for \(S'\to S\),
\(\alpha\in\mathcal O(I_{S'}(\mathcal M))\), and \(f:S'\to F\), by
\[
 \Theta(f\otimes\alpha)
 =\alpha\cdot(\tau_0\circ f)
 =\alpha\cdot(f\circ\rho),
\]
where \(\tau_0:F\to T'_{F/S}(\mathcal M)\) is the zero section and
\(\rho:I_{S'}(\mathcal M)\to S'\) is the structural morphism.

The morphism \(\Theta\) induces
\[
 \theta:F\otimes_{\mathcal O_S}
 \underline{\operatorname{Lie}}(\mathcal O_S/S,\mathcal M)
 \longrightarrow
 \underline{\operatorname{Lie}}'(F/S,\mathcal M).
\]
This follows from the functoriality in \(\mathcal M\) noted after (4),
and can also be seen directly. Indeed,
\[
 \underline{\operatorname{Lie}}'(F/S,\mathcal M)(S')
 =\{\,\varphi\in\operatorname{Hom}_S(I_{S'}(\mathcal M),F)
       \mid \varphi\circ\varepsilon_{\mathcal M}=e\,\},
\]
where \(e:S'\to F\) is the unit section. On the other hand,
\(\underline{\operatorname{Lie}}(\mathcal O_S/S,\mathcal M)(S')\)
equals \(\Gamma(S',\mathcal M)\), the kernel of the augmentation
\[
 \eta:\mathcal O(I_{S'}(\mathcal M))\longrightarrow\mathcal O(S').
\]
It remains to see that if \(f\in F(S')\) and
\(\alpha\in\Gamma(S',\mathcal M)\), then
\(\alpha\cdot(f\circ\rho)\circ\varepsilon_{\mathcal M}=e\).

Apply the dihomomorphism (3) to the \(S\)-morphism
\(\varepsilon_{\mathcal M}:S'\to I_{S'}(\mathcal M)\). The following
diagram commutes.

\SGAThreeDiagram[0.43\textwidth]
  {figures/exp2/exp2_p032_epsilon_action_square.png}
  {Compatibility of scalar action with
  \(F(\varepsilon_{\mathcal M})\).}
  {fig:exp2-epsilon-action-square}
  {Polo--Gille Exp2-14oct24.pdf, component-local p. 32; combined-reader p. 90; printed p. 80}

Thus, for every \(\varphi:I_{S'}(\mathcal M)\to F\),
\[
 (\alpha\cdot\varphi)\circ\varepsilon_{\mathcal M}
 =\eta(\alpha)\cdot(\varphi\circ\varepsilon_{\mathcal M}),
\]
and therefore
\((\alpha\cdot\varphi)\circ\varepsilon_{\mathcal M}=e\) when
\(\eta(\alpha)=0\).

In particular, taking \(\mathcal M=t\mathcal O_S\), one has
\(\underline{\operatorname{Lie}}(\mathcal O_S/S)=t\mathcal O_S\), and
the morphism \(F\to\underline{\operatorname{Lie}}'(F/S)\) is
\(f\mapsto t\cdot(f\circ\rho)\).

\footnote{Editor's note: this remark has been added; it will be useful in
4.6.2.}

\subsection*{Remark 4.3.5.}
Let \(F\) be an \(\mathcal O_S\)-module, set
\[
 E=\underline{\operatorname{End}}_{\mathcal O_S\text{-mod.}}(F),
\]
and denote by \(d_F\) and \(d_E\) the \(\mathcal O_S\)-linear morphisms
given by 4.3.2(5):
\[
\begin{aligned}
 d_F:
 F\otimes_{\mathcal O_S}
 \underline{\operatorname{Lie}}(\mathcal O_S/S,\mathcal M)
 &\longrightarrow
 \underline{\operatorname{Lie}}'(F/S,\mathcal M),\\
 d_E:
 E\otimes_{\mathcal O_S}
 \underline{\operatorname{Lie}}(\mathcal O_S/S,\mathcal M)
 &\longrightarrow
 \underline{\operatorname{Lie}}'(E/S,\mathcal M).
\end{aligned}
\]
Remark~4.3.4 gives the following commutative diagram of
\(\mathcal O_S\)-linear morphisms.

\SGAThreeDiagram[0.76\textwidth]
  {figures/exp2/exp2_p032_differentials_square.png}
  {Compatibility of \(d_E\) and \(d_F\).}
  {fig:exp2-differentials-square}
  {Polo--Gille Exp2-14oct24.pdf, component-local p. 32; combined-reader p. 90; printed p. 80}

The right vertical arrow is the isomorphism \((*)\) of 4.3. Thus, by
that result, if \(F/S\) satisfies \(E\), then \(E/S\) satisfies \(E\), and
\((*)\) is also an isomorphism of \(\mathcal O_S\)-modules when the
right-hand terms are equipped with the module structures obtained from
\(E\).
\footnote{Editor's note: this remark has been added; it will be useful in
4.5 and 4.7.}

% Source page: Exp2 component-local p. 33; combined-reader p. 91; printed p. 81.
\subsection*{Remark 4.4.0.}
In 4.3.2, the morphisms (4) are isomorphisms if and only if the
morphisms (5) are. If these equivalent conditions hold, then \(F/S\)
satisfies \(E\). Indeed, it is enough to verify that
\[
 \underline{\operatorname{Lie}}'(F/S,\mathcal M\oplus\mathcal N)
 \simeq
 \underline{\operatorname{Lie}}'(F/S,\mathcal M)
 \times_S
 \underline{\operatorname{Lie}}'(F/S,\mathcal N).
\]
Under the hypothesis, the horizontal arrows in the following commutative
diagram are isomorphisms.

\SGAThreeDiagram[0.84\textwidth]
  {figures/exp2/exp2_p033_good_module_square.png}
  {The direct-sum comparison used to verify condition \(E\).}
  {fig:exp2-good-module-square}
  {Polo--Gille Exp2-14oct24.pdf, component-local p. 33; combined-reader p. 91; printed p. 81}

The second vertical arrow is therefore an isomorphism as well, which is
exactly condition \(E\) for \(F/S\).
\footnote{Editor's note: in view of the changes made in Definition 4.4
(cf. editor's note (65)), this remark has been inserted here; in the
original it occurred in the proof of 4.4.1.}

\subsection*{Definition 4.4.}
An \(\mathcal O_S\)-module \(F\) is called \emph{good} if the morphisms
\[
 F\otimes_{\mathcal O_S}T_{\mathcal O_S/S}(\mathcal M)
 \longrightarrow T_{F/S}(\mathcal M),
\]
or, equivalently,
\[
 F\otimes_{\mathcal O_S}
 \underline{\operatorname{Lie}}(\mathcal O_S/S,\mathcal M)
 \longrightarrow
 \underline{\operatorname{Lie}}(F/S,\mathcal M),
\]
are isomorphisms of abelian \(S\)-groups (so that \(F/S\) satisfies
(E)), and if moreover they respect the \(\mathcal O_S\)-module
structures obtained from condition \(E\).

\subsection*{Corollary 4.4.1.}
Let \(F\) be an \(\mathcal O_S\)-module, and consider the following
conditions:\footnote{Editor's note: the original proved that (i) implies
(ii) and (iii); the implication (ii) implies (i) has been added.}
\begin{enumerate}
  \item[(i)] \(F\) is a good \(\mathcal O_S\)-module;
  \item[(ii)] \(F/S\) satisfies \(E\), and
    \(d:F\to\underline{\operatorname{Lie}}(F/S)\) is an isomorphism of
    \(\mathcal O_S\)-modules;
  \item[(iii)]
    \(\underline{\operatorname{Lie}}(F/S,\mathcal M)
      =\underline{\operatorname{Lie}}'(F/S,\mathcal M)\).
\end{enumerate}
Then
\[
  \textup{(i)}\Longleftrightarrow\textup{(ii)}
  \Longrightarrow\textup{(iii)}.
\]

\paragraph{Proof.}
The implication
\(\textup{(i)}\Rightarrow\textup{(ii)}\) follows from the definition. To
prove \(\textup{(ii)}\Rightarrow\textup{(i)}\), it remains to show that
the morphisms of
abelian \(S\)-groups, functorial in \(\mathcal M\),
\[
 F\otimes_{\mathcal O_S}
 \underline{\operatorname{Lie}}(\mathcal O_S/S,\mathcal M)
 \xrightarrow{\sim}
 \underline{\operatorname{Lie}}(F/S,\mathcal M)
\]
are isomorphisms of \(\mathcal O_S\)-modules. Since \(F/S\) satisfies
\(E\), both sides carry finite direct sums of copies of \(\mathcal O_S\)
to finite products of abelian \(S\)-groups. We are therefore reduced to
the case \(\mathcal M=\mathcal O_S\), which is precisely the hypothesis.

Finally, \(\textup{(i)}\Rightarrow\textup{(iii)}\) follows from the
definition and from the
fact that the isomorphisms
\[
 F\otimes_{\mathcal O_S}
 \underline{\operatorname{Lie}}(\mathcal O_S/S,\mathcal M)
 \xrightarrow{\sim}
 \underline{\operatorname{Lie}}'(F/S,\mathcal M)
\]
of 4.3.2(5) are \(\mathcal O_S\)-linear.

\subsection*{Examples 4.4.2.}
For every quasi-coherent \(\mathcal O_S\)-module \(\mathcal E\), the
\(\mathcal O_S\)-modules \(\mathbf V(\mathcal E)\) and
\(\mathbf W(\mathcal E)\) defined in I~4.6 are good.

% Source page: Exp2 component-local p. 34; combined-reader p. 92; printed p. 82.
Indeed, for every \(f:S'\to S\), the morphisms
\[
\begin{aligned}
 \mathbf V(\mathcal E)(S')\otimes_{\mathcal O(S')}
   \mathcal O(I_{S'}(\mathcal M))
 &\longrightarrow T_{\mathbf V(\mathcal E)/S}(\mathcal M)(S'),\\
 \mathbf W(\mathcal E)(S')\otimes_{\mathcal O(S')}
   \mathcal O(I_{S'}(\mathcal M))
 &\longrightarrow T_{\mathbf W(\mathcal E)/S}(\mathcal M)(S')
\end{aligned}
\]
correspond respectively to
\[
\begin{aligned}
 &\operatorname{Hom}_{\mathcal O_{S'}\text{-mod.}}
   \bigl(f^*(\mathcal E),\mathcal O_{S'}\bigr)
 \otimes_{\mathcal O(S')}
 \Gamma\bigl(S',D_{\mathcal O_{S'}}(\mathcal M)\bigr)\\
 &\hspace{5em}\longrightarrow
 \operatorname{Hom}_{\mathcal O_{S'}\text{-mod.}}
   \bigl(f^*(\mathcal E),D_{\mathcal O_{S'}}(\mathcal M)\bigr)
 \qquad\textup{(V)}
\\[-2pt]
 &\Gamma(S',f^*(\mathcal E))
 \otimes_{\mathcal O(S')}
 \Gamma\bigl(S',D_{\mathcal O_{S'}}(\mathcal M)\bigr)\\
 &\hspace{5em}\longrightarrow
 \Gamma\bigl(S',f^*(\mathcal E)
   \otimes_{\mathcal O_{S'}}D_{\mathcal O_{S'}}(\mathcal M)\bigr)
 \qquad\textup{(W)}.
\end{aligned}
\]
These are isomorphisms because, as an \(\mathcal O_{S'}\)-module,
\(D_{\mathcal O_{S'}}(\mathcal M)\) is isomorphic to a finite direct
sum of copies of \(\mathcal O_{S'}\).
\footnote{Editor's note: the preceding proof has been added.}

% End of the assigned scope. Proposition 4.5 begins immediately here in
% Exp2-14oct24.pdf, component-local p. 34 (combined-reader p. 92).
